\chapter{Probabilidad: fundamentos y ley de los grandes
números}\label{ch:b3:probability}

En segundo año se construyó la probabilidad sobre espacios numerables;
la teoría de la \omterm{def:b3:measure:measure}{medida} elimina ahora
toda restricción. Un \omterm{def:b3:probability:space}{espacio de
probabilidad} es un espacio medido de
masa total $1$, las \omterm{def:b3:probability:space}{variables
aleatorias} son aplicaciones
\omterm{def:b3:lebesgue:measurable}{medibles}, la
\omterm{def:b3:probability:space}{esperanza} es la integral de Lebesgue
--- y, de golpe, todo el arsenal analítico
(el~\cref{ch:b3:measure,ch:b3:lebesgue,ch:b3:product}) se aplica al
azar. Este capítulo instala el diccionario, construye sucesiones
infinitas de \omterm{def:b3:probability:space}{variables aleatorias}
\omterm{def:b3:probability:independence}{independientes} (en $\intcc01$,
a partir de dígitos binarios: el azar se esconde dentro de la
\omterm{def:b3:measure:lebesgueouter}{medida de Lebesgue}), demuestra
los lemas de Borel--Cantelli y la
\omterm{thm:b3:probability:zeroone}{ley cero--uno} de Kolmogorov, ordena
los modos de convergencia y demuestra la
\omterm{def:b3:probability:space}{ley} de los grandes números --- el
teorema que hace converger las frecuencias hacia las probabilidades y
hace posible la estadística. El problema de fin de semana da la
demostración de Etemadi de la \omterm{def:b3:probability:space}{ley}
fuerte en su forma definitiva $L^1$.

\section{El diccionario}

\begin{definition}\label{def:b3:probability:space}
Un \emph{espacio de probabilidad}\index{espacio de probabilidad} es un
espacio medido
$(\Omega, \mathcal A, \P)$ con $\P(\Omega) =
1$; los elementos de $\mathcal A$ son \emph{sucesos}, y una propiedad se
cumple \emph{casi seguramente} (c.s.) si su suceso tiene probabilidad
$1$. Una \emph{variable aleatoria}\index{variable aleatoria} es una
aplicación \omterm{def:b3:lebesgue:measurable}{medible}
$X \colon \Omega \to \R$ (o $\R^d$: un vector aleatorio); su
\emph{ley}\index{ley de una variable aleatoria} es la
\omterm{def:b3:measure:measure}{medida} de probabilidad imagen $\P_X =
X_*\P$ en $\R$ (el~\cref{exo:b3:product:9}), determinada por la
\emph{función de distribución} $F_X(t) = \P(X \leq t)$
(el~\cref{exo:b3:measure:3}). $X$ tiene
\emph{\omterm{ex:b3:lebesgue:gamma}{densidad}} $f$ si
$\P_X = f\,\dd\lambda$; es \emph{discreta} si $\P_X$ es una combinación
numerable de masas de Dirac. La \emph{esperanza}\index{esperanza} es
\[
\E[X] = \int_\Omega X\,\dd\P
\qquad (X \geq 0 \text{ o } X \in L^1(\P)),
\]
y el \emph{teorema de transferencia} (el~\cref{exo:b3:product:9}) la
calcula en la ley: $\E[g(X)] = \int_\R g\,\dd\P_X$ ---
$= \sum g(x_k)p_k$ en el caso discreto, $= \int
g(x)f(x)\dd x$ en el caso con \omterm{ex:b3:lebesgue:gamma}{densidad}:
las fórmulas de segundo año, ahora teoremas de una sola teoría. La
\emph{varianza} es $\V(X) =
\E[(X - \E X)^2] = \E[X^2] - (\E X)^2$ para $X \in L^2$.
\end{definition}

\begin{example}\label{ex:b3:probability:laws}
Las \omterm{def:b3:probability:space}{leyes} estándar y sus
transformadas destacables: Bernoulli $\mathcal B(p)$, binomial
$\mathcal B(n, p)$, geométrica, Poisson $\mathcal P(\lambda)$
(discretas: las tablas de segundo año siguen siendo válidas); uniforme
en $\intcc01$ (la propia \omterm{def:b3:measure:lebesgueouter}{medida de
Lebesgue}); exponencial $\mathcal E(\lambda)$
(\omterm{ex:b3:lebesgue:gamma}{densidad}
$\lambda\eu^{-\lambda x}\mathbf 1_{x>0}$); la \emph{gaussiana}\index{ley
gaussiana} $\mathcal N(m, \sigma^2)$ de
\omterm{ex:b3:lebesgue:gamma}{densidad}
$\frac1{\sigma\sqrt{2\pi}}\exp\bigl(-\frac{(x -
m)^2}{2\sigma^2}\bigr)$ --- una \omterm{ex:b3:lebesgue:gamma}{densidad}
de probabilidad por el~\cref{pb:b3:lebesgue:1}, de media $m$ y varianza
$\sigma^2$ (momentos gaussianos, el~\cref{exo:b3:product:10}).
\end{example}

\begin{proposition}[Markov y Chebyshev]
\label{prop:b3:probability:markov}
Para $X \geq 0$ y $a > 0$: $\P(X \geq a) \leq \frac{\E
X}{a}$; para $X \in L^2$: $\P\bigl(\abs{X - \E X} \geq
a\bigr) \leq \frac{\V(X)}{a^2}$.\index{desigualdad de
Markov}\index{desigualdad de Chebyshev}
\end{proposition}

\begin{proof}
El~\cref{exo:b3:lebesgue:5}(a); Chebyshev es Markov aplicado a
$(X - \E X)^2$.
\end{proof}

\section{Independencia}

\begin{definition}\label{def:b3:probability:independence}
Las sub-$\sigma$-álgebras $\mathcal A_1, \dots, \mathcal A_n
\subseteq \mathcal A$ son \emph{independientes}\index{independencia} si
$\P(A_1\cap\dots\cap A_n) = \prod\P(A_i)$ para todos $A_i \in
\mathcal A_i$; los sucesos son independientes si lo son las
$\sigma$-álgebras $\{\varnothing, A_i, A_i^c, \Omega\}$; las
\omterm{def:b3:probability:space}{variables aleatorias}
$X_1, \dots, X_n$ lo son si lo son las $\sigma$-álgebras $\sigma(X_i) =
X_i^{-1}(\mathcal B(\R))$. Una familia infinita es independiente si toda
subfamilia finita lo es.
\end{definition}

\begin{theorem}\label{thm:b3:probability:independence}
$X_1, \dots, X_n$ son
\omterm{def:b3:probability:independence}{independientes} si y solo si la
\omterm{def:b3:probability:space}{ley} del vector $(X_1, \dots, X_n)$ es
la \omterm{thm:b3:product:existence}{medida producto}
$\P_{X_1}\otimes\cdots\otimes\P_{X_n}$. En tal caso, para $g_i \geq 0$
(o tales que los productos sean
\omterm{def:b3:lebesgue:l1}{integrables}):
\[
\E\Bigl[\prod_ig_i(X_i)\Bigr] = \prod_i\E[g_i(X_i)],
\]
en particular, $\E[XY] = \E X\,\E Y$ y $\V(X_1 + \dots +
X_n) = \sum\V(X_i)$ para variables $L^2$
\omterm{def:b3:probability:independence}{independientes}.
\end{theorem}

\begin{proof}
Si las $X_i$ son
\omterm{def:b3:probability:independence}{independientes}, las dos
\omterm{def:b3:measure:measure}{medidas} de probabilidad
$\P_{(X_1,\dots,X_n)}$ y $\bigotimes\P_{X_i}$ coinciden en todos los
productos $B_1\times\dots\times B_n$ de borelianos --- un $\pi$-sistema
que genera $\mathcal B(\R^n)$ (la~\cref{prop:b3:product:sections}(b))
---, luego en todas partes (el~\cref{thm:b3:measure:uniqueness}).
Recíprocamente, una \omterm{def:b3:probability:space}{ley} producto
factoriza todos los sucesos $\bigcap_iX_i^{-1}(B_i)$:
\omterm{def:b3:probability:independence}{independencia}. La fórmula de
la \omterm{def:b3:probability:space}{esperanza} es entonces
Tonelli/Fubini (el~\cref{thm:b3:product:tonelli}) a través del teorema
de transferencia; $\E[XY] = \E X\E Y$ es el caso $g_i =
\mathrm{id}$, y desarrollar el cuadrado da la aditividad de las
varianzas (los términos cruzados son $\E[(X_i - \E X_i)(X_j - \E X_j)]
= 0$).
\end{proof}

\begin{theorem}[Existencia de sucesiones independientes]
\label{thm:b3:probability:existence}
En $\bigl(\intcc01, \mathcal L, \lambda\bigr)$ existe una sucesión
$(U_n)_{n\geq1}$ de \omterm{def:b3:probability:space}{variables
aleatorias} \omterm{def:b3:probability:independence}{independientes},
cada una uniforme en $\intcc01$. En consecuencia, para cualesquiera
\omterm{def:b3:probability:space}{leyes} prescritas $(\mu_n)$ en $\R$
existen $(X_n)$ \omterm{def:b3:probability:independence}{independientes}
con $\P_{X_n} = \mu_n$.\index{sucesiones independientes, existencia}
\end{theorem}

\begin{proof}
\emph{Dígitos.} Para $\omega \in \intcc01$, sean $(b_k(\omega))$ sus
dígitos binarios ($\omega = \sum b_k2^{-k}$; elíjase el desarrollo que
no termina en una cola de $1$ --- la ambigüedad afecta solo a un
conjunto numerable, luego nulo). Cada $b_k$ es una
\omterm{def:b3:probability:space}{variable aleatoria} ($\{b_k = 1\}$ es
una unión finita de intervalos diádicos) y el vector $(b_1, \dots, b_m)$
toma cada valor de $\{0,1\}^m$ en un intervalo diádico de longitud
$2^{-m}$: los $b_k$ son Bernoulli$(\frac12)$
\omterm{def:b3:probability:independence}{independientes}.

\emph{Reagrupación.} Pártase $\N^*$ en infinitos conjuntos infinitos
disjuntos $(I_n)$ (por ejemplo, mediante potencias de primos, o
diagonales); sea $(k^n_j)_j$ una enumeración de $I_n$ y póngase
\[
U_n = \sum_{j\geq1} b_{k^n_j}\,2^{-j} .
\]
Cada $U_n$ es uniforme: sus dígitos binarios son bits equilibrados
\omterm{def:b3:probability:independence}{independientes}, de modo que
$\P(U_n \in [l2^{-m}, (l+1)2^{-m})) = 2^{-m}$ para todo intervalo
diádico, y los intervalos diádicos determinan la
\omterm{def:b3:probability:space}{ley}
(el~\cref{thm:b3:measure:uniqueness}). Las $U_n$ son
\omterm{def:b3:probability:independence}{independientes}: son funciones
de bloques disjuntos de la familia
independiente $(b_k)$ ---
formalmente, los sucesos $\{U_n \in D_n\}$ para $D_n$ diádicos dependen
de un número finito de dígitos de conjuntos disjuntos y factorizan; el
argumento del $\pi$-sistema lo eleva a todos los borelianos.

\emph{\omterm{def:b3:probability:space}{Leyes} arbitrarias.} Sea
$G_n(u) = \inf\{t : F_{\mu_n}(t)
\geq u\}$ (la \emph{función cuantil} de la función de distribución
$F_{\mu_n}$); la equivalencia clave $G_n(u) \leq t
\iff u \leq F_{\mu_n}(t)$
(\omterm{def:b3:topology:continuity}{continuidad} por la derecha de $F$,
monotonía) muestra que $X_n = G_n(U_n)$ es
\omterm{def:b3:lebesgue:measurable}{medible} con
$\P(X_n \leq t) = \P(U_n \leq F_{\mu_n}(t)) =
F_{\mu_n}(t)$: \omterm{def:b3:probability:space}{ley} $\mu_n$; la
\omterm{def:b3:probability:independence}{independencia} se hereda
(funciones de variables
\omterm{def:b3:probability:independence}{independientes},
\cref{exo:b3:probability:3}).
\end{proof}

\begin{example}[El problema de los cumpleaños, honestamente]
\label{ex:b3:probability:birthday}
Entre $n$ personas con cumpleaños
\omterm{def:b3:probability:independence}{independientes} y uniformes
sobre $N = 365$ días, la probabilidad de que todos los cumpleaños sean
distintos vale
\[
p_n = \prod_{k=1}^{n-1}\Bigl(1 - \frac kN\Bigr),
\]
por condicionamientos sucesivos (o directamente: los
$N(N-1)\cdots(N - n + 1)$ casos favorables entre los $N^n$ totales, un
argumento de recuento que la fórmula del producto de la
\omterm{def:b3:probability:independence}{independencia} hace riguroso).
Tomando logaritmos y usando $-\ln(1 - x) = x +
O(x^2)$:
\[
\ln p_n = -\frac{n(n-1)}{2N} +
O\Bigl(\frac{n^3}{N^2}\Bigr),
\qquad\text{luego}\qquad
p_n \approx \eu^{-n^2/2N} .
\]
El punto de inflexión $p_n = \frac12$ se sitúa en $n \approx
\sqrt{2N\ln2} \approx 1.18\sqrt N$: para $N = 365$, $n = 23$
($p_{23} = 0.4927$). Dos moralejas. La primera: las colisiones entre $n$
objetos en $N$ casillas aparecen a escala $n \sim \sqrt N$, no
$n \sim N$ --- el \emph{escalado del cumpleaños} que rige las colisiones
de las funciones de dispersión y el coste $\sqrt N$ de los ataques del
cumpleaños en criptografía. La segunda: el cálculo es una plantilla: los
$\binom n2$ sucesos de colisión por pares no son
\omterm{def:b3:probability:independence}{independientes} y, sin embargo,
la respuesta se comporta como si lo fueran ($\eu^{-\binom n2/N}$ es
exactamente la heurística de pares \omterm{def:b3:probability:independence}{independientes}) --- una primera
instancia de la aproximación de Poisson, que el problema de fin de
semana del~\cref{ch:b3:clt} hace rigurosa (desigualdad de Le Cam).
\end{example}

\section{Borel--Cantelli y la ley cero--uno}

\begin{theorem}[Borel--Cantelli]
\label{thm:b3:probability:borelcantelli}
Sean $(A_n)$ sucesos y $\limsup A_n = \bigcap_N
\bigcup_{n\geq N}A_n$ («$A_n$ ocurre infinitas veces»).\index{lemas de
Borel--Cantelli}
\begin{enumerate}
  \item Si $\sum\P(A_n) < \infty$, entonces $\P(\limsup A_n) =
        0$. \item Si $\sum\P(A_n) = \infty$ \emph{y los $A_n$ son
  \omterm{def:b3:probability:independence}{independientes}}, entonces
  $\P(\limsup A_n) = 1$.
\end{enumerate}
\end{theorem}

\begin{proof}
(1) es el~\cref{exo:b3:measure:4}. (2): para $N \leq M$, la
\omterm{def:b3:probability:independence}{independencia} de los
complementarios (el~\cref{exo:b3:probability:3}) da
\[
\P\Bigl(\bigcap_{n=N}^{M}A_n^c\Bigr) = \prod_{n=N}^M\bigl(1
- \P(A_n)\bigr) \leq
\exp\Bigl(-\sum_{n=N}^M\P(A_n)\Bigr) \xrightarrow[M \to
\infty]{} 0
\]
($1 - x \leq \eu^{-x}$; la serie diverge). Así,
$\P\bigl(\bigcup_{n\geq N}A_n\bigr) = 1$ para todo $N$, y la
intersección decreciente en $N$ sigue teniendo probabilidad $1$
(\omterm{def:b3:topology:continuity}{continuidad} por arriba,
la~\cref{prop:b3:measure:basics}).
\end{proof}

\begin{theorem}[Ley cero--uno de Kolmogorov]
\label{thm:b3:probability:zeroone}
Sean $(X_n)$ \omterm{def:b3:probability:independence}{independientes} y
$\mathcal T =
\bigcap_N\sigma(X_N, X_{N+1}, \dots)$ la \emph{$\sigma$-álgebra de
cola}\index{$\sigma$-álgebra de cola} (sucesos insensibles a cualquier
número finito de las $X_n$: convergencia de $\sum X_n$, de
$\frac{S_n}n$, valores de $\limsup$, \dots). Entonces todo
$T \in \mathcal T$ cumple $\P(T) \in \{0,
1\}$.\index{ley cero--uno}
\end{theorem}

\begin{proof}
Fíjese $N$. Las $\sigma$-álgebras $\sigma(X_1, \dots, X_N)$ y
$\sigma(X_{N+1}, \dots)$ son
\omterm{def:b3:probability:independence}{independientes}: los sucesos
que dependen de bloques disjuntos factorizan sobre los $\pi$-sistemas
generadores (cilindros $\bigcap_{i\leq N}\{X_i \in B_i\}$ y,
respectivamente, condiciones finitas sobre las variables posteriores), y
Dynkin (el~\cref{thm:b3:measure:dynkin}, aplicado dos veces, un lado
cada vez) extiende la factorización. Un suceso de cola $T$ está en
$\sigma(X_{N+1}, \dots)$ para todo $N$: $T$ es
independiente de cada
$\sigma(X_1, \dots, X_N)$, luego de la $\sigma$-álgebra que estas
generan, $\sigma(X_1, X_2, \dots)$ (Dynkin una vez más: la unión de las
$\sigma(X_1,\dots,X_N)$ es un $\pi$-sistema que la genera). Pero también
$T \in \sigma(X_1, X_2, \dots)$: $T$ es
independiente \emph{de sí
mismo}, $\P(T) = \P(T\cap T) =
\P(T)^2$: $\P(T) \in \{0, 1\}$.
\end{proof}

\section{Modos de convergencia}

\begin{definition}\label{def:b3:probability:modes}
$X_n \to X$ \emph{casi seguramente} si $\P(X_n \to X) = 1$; \emph{en
probabilidad} si $\P(\abs{X_n - X} \geq \varepsilon)
\to 0$ para todo $\varepsilon > 0$; \emph{en $L^p$} si
$\E\abs{X_n - X}^p \to 0$.\index{convergencia de variables aleatorias}
\end{definition}

\begin{proposition}\label{prop:b3:probability:modes}
(a) la convergencia c.s.\ implica la convergencia en probabilidad; (b)
la convergencia en $L^p$ implica la convergencia en probabilidad; (c) la
convergencia en probabilidad implica la convergencia c.s.\ \emph{a lo
largo de una subsucesión}; (d) ninguna otra implicación es cierta en
general.
\end{proposition}

\begin{proof}
(a) $\P(\abs{X_n - X} \geq \varepsilon) \leq
\P\bigl(\sup_{m\geq n}\abs{X_m - X} \geq \varepsilon\bigr)
\downarrow \P\bigl(\limsup\{\abs{X_m - X} \geq
\varepsilon\}\bigr) = 0$ bajo convergencia c.s.\
(\omterm{def:b3:topology:continuity}{continuidad} por arriba; el suceso
del límite superior excluye la convergencia). (b) Markov:
$\P(\abs{X_n - X} \geq \varepsilon) \leq
\varepsilon^{-p}\,\E\abs{X_n - X}^p$. (c) Tómese $n_k$ con
$\P(\abs{X_{n_k} - X} \geq 2^{-k}) \leq 2^{-k}$;
Borel--Cantelli (1) hace que $\abs{X_{n_k} - X} < 2^{-k}$ a partir de
cierto índice, c.s. (d) La máquina de escribir (el~\cref{exo:b3:lp:3})
en $(\intcc01, \lambda)$ converge en $L^1$ y en probabilidad, pero en
ningún punto; $n\mathbf 1_{\intoo0{1/n}} \to 0$ c.s.\ pero no en $L^1$;
los detalles y los contraejemplos restantes están en
el~\cref{exo:b3:probability:6}.
\end{proof}

\section{La ley de los grandes números}

En todo lo que sigue, $(X_n)$ son
\omterm{def:b3:probability:independence}{independientes} con la misma
\omterm{def:b3:probability:space}{ley}
(\emph{i.i.d.}), $S_n = X_1 + \dots + X_n$.

\begin{theorem}[Ley débil de los grandes números]
\label{thm:b3:probability:wlln}
Si $X_1 \in L^2$, con $m = \E X_1$:
\[
\P\Bigl(\Bigl|\frac{S_n}{n} - m\Bigr| \geq
\varepsilon\Bigr) \leq
\frac{\V(X_1)}{n\,\varepsilon^2}
\xrightarrow[n\to\infty]{} 0 :
\]
$\frac{S_n}n \to m$ en probabilidad (y en $L^2$).\index{ley de los
grandes números (débil)}
\end{theorem}

\begin{proof}
$\E\frac{S_n}n = m$ y $\V\bigl(\frac{S_n}n\bigr) =
\frac{n\V(X_1)}{n^2}$ (el~\cref{thm:b3:probability:independence});
Chebyshev.
\end{proof}

\begin{theorem}[Ley fuerte de los grandes números]
\label{thm:b3:probability:slln}
Si $X_1 \in L^1$, entonces\index{ley de los grandes números (fuerte)}
\[
\frac{S_n}{n} \xrightarrow[n\to\infty]{\text{c.s.}} \E[X_1].
\]
La demostramos aquí bajo la hipótesis más fuerte $X_1 \in
L^4$; el caso general ($L^1$: la demostración de Etemadi) es el problema
de fin de semana.
\end{theorem}

\begin{proof}[Demostración bajo $\E X_1^4 < \infty$]
Centrando ($X_i \mapsto X_i - m$), supóngase $m = 0$. Desarróllese:
\[
\E[S_n^4] = \sum_{i,j,k,l}\E[X_iX_jX_kX_l]
= n\,\E[X_1^4] + 3n(n-1)\,\bigl(\E[X_1^2]\bigr)^2 \leq
C\,n^2 ,
\]
puesto que la \omterm{def:b3:probability:independence}{independencia} y
el centrado matan todo término que contenga un factor aislado
($\E[X_iX_jX_kX_l] =
\E[X_i]\E[\cdots] = 0$ salvo que los índices se emparejen: los únicos
supervivientes son los $n$ términos $i=j=k=l$ y los $3n(n-1)$ términos
con dos pares distintos). Markov:
\[
\P\Bigl(\Bigl|\frac{S_n}n\Bigr| \geq \varepsilon\Bigr)
= \P\bigl(S_n^4 \geq n^4\varepsilon^4\bigr)
\leq \frac{Cn^2}{n^4\varepsilon^4} =
\frac{C}{\varepsilon^4n^2},
\]
sumable: Borel--Cantelli (1) da, para cada racional $\varepsilon$, que
$\abs{S_n/n} < \varepsilon$ a partir de cierto índice, c.s.;
intersecando en $\varepsilon \in \Q_+^*$ (una cantidad numerable de
sucesos de probabilidad $1$): $S_n/n \to 0$ c.s.
\end{proof}

\begin{example}[Lo que compra la ley fuerte]
\label{ex:b3:probability:sllnapps}
(a) \emph{Frecuencias}: para lanzamientos de moneda i.i.d., la
frecuencia observada de caras converge c.s.\ a $p$ --- la justificación
empírica de la propia probabilidad. (b) \emph{Monte Carlo}\index{método
de Monte Carlo}: para $g \in
L^1(\intcc01)$ y $(U_n)$ uniformes i.i.d.,
(\cref{thm:b3:probability:existence}),
$\frac1n\sum_{k\leq n}g(U_k) \to \int_0^1g$ c.s.: integrales por
muestreo, en cualquier dimensión, a la velocidad
independiente de la dimensión
$\sim n^{-1/2}$ que precisa el~\cref{ch:b3:clt}. (c) \emph{Números
normales}: casi todo número real tiene, en su desarrollo binario,
frecuencia asintótica $\frac12$ de unos (aplíquese la
\omterm{def:b3:probability:space}{ley} fuerte a las variables de dígitos
del~\cref{thm:b3:probability:existence}) --- el teorema de Borel, un
enunciado sobre los números de \emph{todos} los días demostrado por la
\omterm{def:b3:measure:measure}{medida}: el~\cref{pb:b3:probability:1}
lo \omterm{def:b3:complete:complete}{completa} en todas las bases.
\end{example}

\begin{method}\label{met:b3:probability:toolkit}
El orden de trabajo para los enunciados asintóticos sobre sucesiones
aleatorias: (1) \emph{¿es el suceso un suceso de cola?} Entonces su
probabilidad vale $0$ o $1$ (el~\cref{thm:b3:probability:zeroone}) y
solo hay que decidir cuál. (2) \emph{Para demostrar enunciados c.s.}:
Borel--Cantelli --- probabilidades sumables para los sucesos «malos»,
mediante cotas de tipo Markov o Chebyshev sobre los momentos que
existan; la \omterm{def:b3:probability:independence}{independencia} solo
hace falta en el sentido recíproco. (3) \emph{Subsucesión más sándwich}:
demuéstrese la convergencia a lo largo de una subsucesión manejable y
contrólese la oscilación intermedia por monotonía o desigualdades
maximales --- el esqueleto de la demostración de Etemadi. (4) Para los
límites en distribución, espérese al~\cref{ch:b3:clt}.
\end{method}

\section{Ejercicios}

\begin{exercise}[$\star$]\label{exo:b3:probability:1}
(a) Sea $X$ de función de distribución $F$
\omterm{def:b3:topology:continuity}{continua} y estrictamente creciente.
Demuéstrese que $F(X)$ es uniforme en $\intcc01$ y que $G(U) \sim F$
para $U$ uniforme, $G = F^{-1}$: simulación por inversión. (b)
Calcúlense la función de distribución y la
\omterm{ex:b3:lebesgue:gamma}{densidad} de $X^2$ para $X$ uniforme en
$\intcc{-1}1$, y de $-\frac1\lambda\ln
U$ para $U$ uniforme en $\intoo01$.
\end{exercise}

\begin{exercise}[$\star$]\label{exo:b3:probability:2}
(a) Calcúlense la media y la varianza de las
\omterm{def:b3:probability:space}{leyes} de Poisson $\mathcal
P(\lambda)$ y geométrica mediante el teorema de transferencia. (b)
Demuéstrese que una \omterm{def:b3:probability:space}{variable
aleatoria} positiva $T$ con $\P(T > t)
> 0$ para todo $t$ cumple la propiedad de \emph{falta de memoria}
$\P(T > t + s \mid
T > t) = \P(T > s)$ para todos $s, t \geq 0$ si y solo si $T$ es
exponencial. \emph{(La función de supervivencia satisface la ecuación
funcional de Cauchy; la monotonía sustituye a la
\omterm{def:b3:topology:continuity}{continuidad}.)}
\end{exercise}

\begin{exercise}[$\star\star$]\label{exo:b3:probability:3}
(a) Demuéstrese que si $X_1, \dots, X_n$ son
\omterm{def:b3:probability:independence}{independientes} y $f_i$ son
funciones borelianas, las $f_i(X_i)$ son
\omterm{def:b3:probability:independence}{independientes}. (b)
Demuéstrese que unos sucesos $A_1, \dots, A_n$ son
\omterm{def:b3:probability:independence}{independientes} si y solo si lo
son sus complementarios, si y solo si los indicadores $\mathbf 1_{A_i}$
son \omterm{def:b3:probability:space}{variables aleatorias}
\omterm{def:b3:probability:independence}{independientes}. (c) (La
\omterm{def:b3:probability:independence}{independencia} por pares es más débil) Dos monedas equilibradas: $A = $
la primera sale cara, $B = $ la segunda sale cara, $C = $ ambas
coinciden. Demuéstrese que $A, B, C$ son
\omterm{def:b3:probability:independence}{independientes} dos a dos pero
no \omterm{def:b3:probability:independence}{independientes}.
\end{exercise}

\begin{exercise}[$\star\star$]\label{exo:b3:probability:4}
(a) (El mono infinito) Una sucesión i.i.d.\ de pulsaciones uniformes
sobre un alfabeto finito contiene c.s.\ todo texto finito infinitas
veces: demuéstrese con Borel--Cantelli (2) sobre bloques disjuntos. (b)
(Rachas) Para bits equilibrados i.i.d., sea $R_n$ la longitud de la
racha de unos que empieza en la posición $n$. Demuéstrese que, c.s.,
$R_n \geq (1+\varepsilon)\log_2n$ un número finito de veces, y $R_n
\geq \log_2 n$ infinitas veces \emph{(ambas mitades de Borel--Cantelli;
para la segunda, pásese a bloques disjuntos para ganar
\omterm{def:b3:probability:independence}{independencia})}: la racha más
larga entre los $n$ primeros dígitos crece como $\log_2n$.
\end{exercise}

\begin{exercise}[$\star\star$]\label{exo:b3:probability:5}
Sean $(X_n)$ \omterm{def:b3:probability:independence}{independientes}.
(a) Demuéstrese que el radio de convergencia de $\sum X_n z^n$ es una
constante c.s.\ (posiblemente $0$ o $\infty$). (b) Demuéstrense
$\P(\sum X_n \text{ converge}) \in \{0, 1\}$ y
$\P(S_n/n \to m) \in \{0,1\}$. (c) Dese un suceso relativo a $(X_n)$ que
\emph{no} sea un suceso de cola, y compruébese que la
\omterm{thm:b3:probability:zeroone}{ley cero--uno} puede fallar para él.
\end{exercise}

\begin{exercise}[$\star\star$]\label{exo:b3:probability:6}
En $(\intcc01, \lambda)$, exhíbanse --- con demostración ---
\omterm{def:b3:probability:space}{variables aleatorias} tales que: (a)
$X_n \to 0$ en probabilidad y en todo $L^p$, pero en ningún punto c.s.;
(b) $X_n \to 0$ c.s.\ pero en ningún $L^p$; (c) $X_n \to 0$ en $L^1$
pero no en $L^2$; (d) y demuéstrese: si $X_n \to X$ en probabilidad y
$\abs{X_n}
\leq Y \in L^1$, entonces $X_n \to X$ en $L^1$ \emph{(subsucesiones,
convergencia dominada y el truco de la subsubsucesión)}.
\end{exercise}

\begin{exercise}[$\star\star$]\label{exo:b3:probability:7}
Una encuesta de opinión estima una proporción desconocida $p$ mediante
la frecuencia empírica $\hat p_n$ de $n$ extracciones
\omterm{def:b3:probability:independence}{independientes}. (a) Chebyshev:
demuéstrese $\P(\abs{\hat p_n - p} \geq \varepsilon)
\leq \frac1{4n\varepsilon^2}$ (úsese $p(1-p) \leq \frac14$). (b)
¿Cuántas extracciones garantizan un error $\leq 3\%$ con probabilidad
$\geq 95\%$ según esta cota? (La respuesta verdadera, vía
el~\cref{ch:b3:clt}, es de unas $1070$: Chebyshev es honesta pero
burda.)
\end{exercise}

\begin{exercise}[$\star\star\star$]\label{exo:b3:probability:8}
(Bernstein) Para $f \in \mathcal C(\intcc01)$, defínase el polinomio de
Bernstein $B_nf(x) =
\sum_{k=0}^n\binom nkx^k(1-x)^{n-k}f\bigl(\frac kn\bigr)$. (a)
Reconózcase $B_nf(x) = \E\bigl[f\bigl(\frac
{S_n}n\bigr)\bigr]$ para $S_n$ binomial $\mathcal B(n, x)$. (b)
Demuéstrese $B_nf \to f$ \emph{uniformemente} en $\intcc01$: sepárese
según $\{\abs{\frac{S_n}n - x} \leq \delta\}$ y su complementario,
usando la \omterm{def:b3:topology:continuity}{continuidad} uniforme y
Chebyshev con la cota uniforme $\V(\frac{S_n}n) \leq \frac1{4n}$. (c)
Conclúyase: una segunda demostración, probabilística, del teorema de
aproximación de Weierstrass (el~\cref{cor:b3:complete:weierstrass}), con
la velocidad explícita
$\norm{B_nf - f}_\infty \leq \frac32\,\omega_f(n^{-1/2})$
para el módulo de \omterm{def:b3:topology:continuity}{continuidad}
$\omega_f$ --- demuéstrese al menos la forma $O(\omega_f(n^{-1/2}))$.
\end{exercise}

\begin{exercise}[$\star\star\star$]\label{exo:b3:probability:9}
(Coleccionista de cupones) Se extraen con reposición, uniformemente,
cromos de $n$ tipos; sea $T_n$ el número de extracciones hasta ver todos
los tipos. (a) Escríbase $T_n = \sum_{k=1}^{n}\tau_k$ con $\tau_k$
geométrica de parámetro $\frac{n - k + 1}n$ y las $\tau_k$
\omterm{def:b3:probability:independence}{independientes}, y dedúzcanse
$\E T_n = n\,H_n \sim n\ln n$ ($H_n$ el número armónico) y $\V(T_n) \leq
\frac{\pi^2}6n^2$. (b) Chebyshev: $\frac{T_n}{n\ln n} \to 1$ en
probabilidad. (c) Afínese con Borel--Cantelli: demuéstrese directamente
$\P(T_n >
\beta n\ln n) \leq n^{1 - \beta}$ para $\beta > 1$ \emph{(cota de la
unión sobre el suceso de que falte algún tipo tras $\beta n\ln n$
extracciones, usando $1 - x \leq \eu^{-x}$)}, y dedúzcase que, a lo
largo de $n = 2^m$, c.s.\ $T_n \leq \beta n\ln
n$ a partir de cierto índice, para todo $\beta > 2$.
\end{exercise}

\begin{exercise}[$\star\star$]\label{exo:b3:probability:10}
Usando la construcción por dígitos
(\cref{thm:b3:probability:existence}):
(a) verifíquese por cálculo directo que $U = \sum b_{2k}2^{-k}$ (los
dígitos de índice par de una $\omega$ uniforme) es uniforme e
independiente de
$V = \sum b_{2k-1}2^{-k}$; (b) dedúzcase una biyección
\omterm{def:b3:lebesgue:measurable}{medible} salvo conjuntos nulos entre
$\intcc01$ y $\intcc01^2$ que conserve la
\omterm{def:b3:measure:measure}{medida}, y coméntese: un número
aleatorio uniforme contiene dos (y una infinidad numerable)
\omterm{def:b3:probability:independence}{independientes} --- compárese
con la curva de Peano (el~\cref{pb:b3:topology:1}), que lograba la
sobreyectividad, pero ni la conservación de la \omterm{def:b3:measure:measure}{medida} ni la
inyectividad.
\end{exercise}

\begin{exercise}[$\star\star$]\label{exo:b3:probability:11}
(Récords) Sean $(X_n)_{n\geq1}$ i.i.d.\ de función de distribución
\omterm{def:b3:topology:continuity}{continua}, y dígase que hay un
\emph{récord} en el instante $n$ si $X_n > \max(X_1, \dots, X_{n-1})$
(el instante $1$ es un récord). Sea $R_n$ el indicador de récord. (a)
Demuéstrese $\P(R_n = 1) = \frac1n$ \emph{(por simetría, cada una de las
$n!$ ordenaciones de $X_1, \dots, X_n$ es igualmente probable y los
empates tienen probabilidad $0$)}. (b) Demuéstrese que los $R_n$ son
\emph{\omterm{def:b3:probability:independence}{independientes}}
\emph{(cuéntense las ordenaciones compatibles con posiciones de récord
prescritas, o argúyase que el orden relativo de $X_1, \dots, X_{n-1}$ es
independiente del rango de
$X_n$ entre ellas)}. (c) Dedúzcase de Borel--Cantelli
(el~\cref{thm:b3:probability:borelcantelli}, ambas mitades) que hay
infinitos récords c.s., pero que los récords en instantes consecutivos
$n, n+1$ ocurren infinitas veces con probabilidad --- ¡decídase cuál!
--- y calcúlese
$\sum_n\P(R_n = 1, R_{n+1} = 1)$.
\end{exercise}

\begin{exercise}[$\star\star$]\label{exo:b3:probability:12}
(La racha más larga de caras) Lánzese una moneda equilibrada infinitas
veces y sea $L_n$ la longitud de la racha más larga de caras
consecutivas entre los $n$ primeros lanzamientos. (a) Demuéstrese que,
para todo $\varepsilon > 0$, c.s.\ $L_n \leq
(1 + \varepsilon)\log_2n$ a partir de cierto índice \emph{(la
probabilidad de que alguna racha de longitud $\ell$ empiece entre los
$n$ primeros lanzamientos es a lo sumo $n2^{-\ell}$; Borel--Cantelli a
lo largo de $n =
2^k$)}. (b) Demuéstrese que, c.s., $L_n \geq (1 - \varepsilon)\log_2n$ a
partir de cierto índice \emph{(pártanse los $n$ primeros lanzamientos en
$\lfloor n/\ell\rfloor$ bloques disjuntos de longitud $\ell =
\lceil(1 - \varepsilon)\log_2n\rceil$; los bloques son
\omterm{def:b3:probability:independence}{independientes}, cada uno todo
caras con probabilidad $2^{-\ell}$, y la probabilidad de que ninguno sea
todo caras es a lo sumo $\exp(-n2^{-\ell}/\ell)$; súmese de nuevo a lo
largo de $n = 2^k$)}. (c) Conclúyase $\frac{L_n}{\log_2n} \to 1$ c.s.:
en un millón de lanzamientos equilibrados cabe esperar una racha de unas
$20$ caras --- y un conjunto de datos sin ella es probablemente
inventado.
\end{exercise}

\section{Problema: la demostración de Etemadi de la ley fuerte}

\begin{problem}[{Problema de fin de semana --- la ley fuerte de los
grandes números para variables i.i.d.\ integrables}]
\label{pb:b3:probability:1}
La \omterm{def:b3:probability:space}{ley} fuerte de Kolmogorov ---
$\frac{S_n}n \to \E X_1$ c.s.\ para $X_n \in L^1$ i.i.d.\ --- tuvo
durante mucho tiempo solo demostraciones intrincadas; en 1981,
N.~Etemadi encontró una de una economía asombrosa, que no usa nada más
allá de este capítulo (y que incluso debilita la
\omterm{def:b3:probability:independence}{independencia} a la
\omterm{def:b3:probability:independence}{independencia} dos a dos). La
seguimos. Sean $(X_n)$
\omterm{def:b3:probability:independence}{independientes} dos a dos,
idénticamente distribuidas e \omterm{def:b3:lebesgue:l1}{integrables};
$m = \E X_1$, $S_n = X_1 + \dots + X_n$.

\medskip
\noindent\textbf{Parte I --- Reducciones.}
\begin{enumerate}
  \item Demostrar que basta tratar $X_n \geq 0$ \emph{(pártase
  $X_n = X_n^+ - X_n^-$: compruébese que las dos mitades vuelven a ser
  i.i.d.\ \omterm{def:b3:probability:independence}{independientes} dos a
  dos e \omterm{def:b3:lebesgue:l1}{integrables})}. Supóngase en
  adelante $X_n \geq 0$. \item (Truncamiento) Sean
  $Y_n = X_n\,\mathbf 1_{X_n \leq n}$ y $S_n^* = Y_1 + \dots + Y_n$.
  Demostrar
        \[
        \sum_{n\geq1}\P(X_n \neq Y_n) =
        \sum_{n\geq1}\P(X_1 > n) \leq \E[X_1] < \infty
        \]
        (el~\cref{exo:b3:product:3}) y dedúzcase, vía Borel--Cantelli,
        que $\frac{S_n - S_n^*}{n} \to 0$ c.s.: basta demostrar
        $\frac{S^*_n}n \to m$
        a.s.
  \item Demostrar $\E Y_n = \E\bigl[X_1\mathbf 1_{X_1\leq
        n}\bigr] \to m$ (convergencia monótona) y, por tanto,
  $\frac1n\sum_{k\leq n}\E Y_k \to m$ (Cesàro): basta demostrar
  $\frac{S_n^* - \E S_n^*}{n} \to 0$
        a.s.
\end{enumerate}

\medskip
\noindent\textbf{Parte II --- La estimación de la varianza.}
\begin{enumerate}[resume]
  \item Demostrar
        \[
        \V(Y_n) \leq \E[Y_n^2] = \E\bigl[X_1^2\,\mathbf
        1_{X_1 \leq n}\bigr]
        \]
        y, usando la fórmula de las capas
        (la~\cref{prop:b3:product:layercake}), la cota clave
        \[
        \sum_{n\geq1}\frac{\V(Y_n)}{n^2}
        \leq \sum_{n\geq1}\frac1{n^2}\,
        \E\bigl[X_1^2\mathbf 1_{X_1\leq n}\bigr]
        \leq C\,\E[X_1] < \infty
        \]
        \emph{(intercámbiense la suma y la
        \omterm{def:b3:probability:space}{esperanza} --- Tonelli para
        series --- y acótese $\sum_{n \geq
        x}\frac1{n^2} \leq \frac2{\max(x,1)}$ para la estimación
        \omterm{def:b3:topology:interior}{interior} $x^2\sum_{n\geq x}n^{-2} \leq 2x$)}.
\end{enumerate}

\medskip
\noindent\textbf{Parte III --- Convergencia a lo largo de subsucesiones
geométricas.} Fíjese $\alpha > 1$ y sea $k_j =
\lfloor\alpha^j\rfloor$.
\begin{enumerate}[resume]
  \item Usando la
  \omterm{def:b3:probability:independence}{independencia} dos a dos (las
  varianzas se suman, el~\cref{thm:b3:probability:independence} ---
  compruébese que la aditividad de las varianzas solo necesita la
  \omterm{def:b3:probability:independence}{independencia} dos a dos) y
  Chebyshev, demuéstrese, para todo
        $\varepsilon > 0$:
        \[
        \sum_{j\geq1}\P\Bigl(\Bigl|
        \frac{S^*_{k_j} - \E S^*_{k_j}}{k_j}\Bigr| \geq
        \varepsilon\Bigr)
        \leq
        \frac1{\varepsilon^2}\sum_{j\geq1}\frac1{k_j^2}
        \sum_{n\leq k_j}\V(Y_n)
        = \frac1{\varepsilon^2}\sum_{n\geq1}\V(Y_n)
        \sum_{j\,:\,k_j\geq n}\frac1{k_j^2} .
        \]
  \item Demostrar $\sum_{j : k_j \geq n}k_j^{-2} \leq
        \frac{C_\alpha}{n^2}$ \emph{(serie geométrica; atención a la
  parte entera: $k_j \geq \frac{\alpha^j}2$ para el cuidado de tipo
  $\alpha^j \geq 2$)}, y conclúyase con la pregunta 4 y Borel--Cantelli:
        \[
        \frac{S^*_{k_j} - \E S^*_{k_j}}{k_j}
        \xrightarrow[j\to\infty]{\text{c.s.}} 0,
        \qquad\text{luego}\qquad
        \frac{S^*_{k_j}}{k_j} \to m \ \text{c.s.}
        \]
\end{enumerate}

\medskip
\noindent\textbf{Parte IV --- Sándwich y conclusión.}
\begin{enumerate}[resume]
  \item Para $k_j \leq n \leq k_{j+1}$, úsese la monotonía de $S^*_n$
  (¡sumandos no negativos!) para demostrar
        \[
        \frac{k_j}{k_{j+1}}\,\frac{S^*_{k_j}}{k_j}
        \;\leq\; \frac{S^*_n}{n} \;\leq\;
        \frac{k_{j+1}}{k_j}\,\frac{S^*_{k_{j+1}}}{k_{j+1}},
        \]
        y dedúzcase, c.s.:
        \[
        \frac m\alpha \leq \liminf\frac{S^*_n}n \leq
        \limsup\frac{S^*_n}n \leq \alpha\,m .
        \]
  \item Hágase $\alpha \downarrow 1$ a lo largo de una sucesión y
  conclúyase $\frac{S_n^*}n \to m$ c.s.\ y, por tanto (Parte I), la
  \emph{\omterm{def:b3:probability:space}{ley} fuerte de los grandes
  números}:
        \[
        \boxed{\ \frac{S_n}{n}
        \xrightarrow[n\to\infty]{\text{c.s.}} \E[X_1].\ }
        \]
  \item ¿Dónde bastó exactamente la
  \omterm{def:b3:probability:independence}{independencia} dos a dos (en
  lugar de la \omterm{def:b3:probability:independence}{independencia}
  \omterm{def:b3:complete:complete}{completa})? Enumérense los tres lugares donde se invocaron hipótesis de
  tipo \omterm{def:b3:probability:independence}{independencia}.
\end{enumerate}

\medskip
\noindent\textbf{Parte V --- Dividendos.}
\begin{enumerate}[resume]
  \item (Números normales de Borel) Demostrar que $\lambda$-casi todo
  $x \in \intcc01$ es \emph{normal en toda base} $b \geq 2$: cada dígito
  $0, \dots, b-1$ aparece con frecuencia asintótica $\frac1b$
  \emph{(fíjense $b$ y un dígito, aplíquese la
  \omterm{def:b3:probability:space}{ley} fuerte a las variables
  indicadoras --- justifíquese que los dígitos en base $b$ de una
  variable uniforme son i.i.d.\ uniformes en $\{0,\dots,b-1\}$, como en
  el~\cref{thm:b3:probability:existence} --- e interséquense después los
  sucesos de probabilidad uno, en cantidad numerable)}. Exhíbase un
  número explícito no normal y reflexiónese: el teorema afirma la
  normalidad de casi todos los números y, sin embargo, demostrar la
  normalidad de $\sqrt2$ o de $\pi$ sigue abierto. \item
  (\omterm{ex:b3:probability:sllnapps}{Monte Carlo}, garantizado)
  Justifíquese por completo el método
  del~\cref{ex:b3:probability:sllnapps}(b) para $g \in L^1(\intcc01^d)$:
  constrúyase la muestra uniforme i.i.d.\ en $\intcc01^d$ a partir
  del~\cref{thm:b3:probability:existence} y
  del~\cref{exo:b3:probability:10}, y enúnciese qué entrega la
  \omterm{def:b3:probability:space}{ley} fuerte.
\end{enumerate}

\medskip
\noindent\textbf{Parte VI --- Lo que compra la
\omterm{def:b3:probability:independence}{independencia} \omterm{def:b3:complete:complete}{completa}:
desigualdades maximales y series aleatorias.} Etemadi gasta solo
\omterm{def:b3:probability:independence}{independencia} dos a dos; las
partes restantes explotan la versión \omterm{def:b3:complete:complete}{completa} (mutua). Sean $(Z_n)$
variables \omterm{def:b3:probability:independence}{independientes}
centradas de $L^2$ y $S_k = Z_1 + \dots + Z_k$ (una notación nueva, sin
relación con las $X_n$ anteriores).
\begin{enumerate}[resume]
  \item (Desigualdad maximal de Kolmogorov) Para $\varepsilon
        > 0$, demuéstrese
        \[
        \P\Bigl(\max_{1\leq k\leq n}\abs{S_k} \geq
        \varepsilon\Bigr) \;\leq\;
        \frac1{\varepsilon^2}\sum_{k=1}^n\V(Z_k) :
        \]
        El precio de Chebyshev compra el máximo \emph{(divídase el
        suceso según el primer índice $k$ con
        $\abs{S_k} \geq \varepsilon$; en ese trozo escríbase
        $S_n^2 \geq S_k^2 + 2S_k(S_n - S_k)$ y úsese la
        \omterm{def:b3:probability:independence}{independencia} de las
        coaliciones $(Z_1, \dots, Z_k)$ y $(Z_{k+1}, \dots, Z_n)$,
        el~\cref{thm:b3:probability:independence})}. Señálese el paso en
        que la \omterm{def:b3:probability:independence}{independencia}
        dos a dos ya no bastaría. \item (Teorema de la serie única de
        Khinchin--Kolmogorov) Dedúzcase: si $\sum_n\V(Z_n) < \infty$,
        entonces $\sum_nZ_n$ converge casi seguramente
        \emph{(demuéstrese que, c.s., las sumas parciales forman una
        sucesión de Cauchy: hágase $m \to
        \infty$ en la desigualdad maximal aplicada a
        $Z_{N+1}, \dots, Z_{N+m}$, y después $N \to
        \infty$)}. \item (Series de Rademacher) Sean $(\varepsilon_n)$
        signos i.i.d., $\P(\varepsilon_n = \pm1) =
        \frac12$ (el~\cref{thm:b3:probability:existence}), y sean
        $(x_n)$ números reales. Demuéstrese que $\sum_nx_n\varepsilon_n$
        converge c.s.\ en cuanto $\sum_nx_n^2 < \infty$; demuéstrese
        también que, sean cuales sean los $(x_n)$, la probabilidad de
        que $\sum_nx_n\varepsilon_n$ converja vale $0$ o $1$
        (\cref{thm:b3:probability:zeroone}).
  \item El recíproco, elementalmente. Póngase $T_n = \sum_{k\leq
        n}x_k\varepsilon_k$ y $s_n^2 = \sum_{k\leq
        n}x_k^2$, y supóngase $s_n \to \infty$. (a) Demuéstrese la
  \emph{desigualdad de Paley--Zygmund}: para $Z \geq 0$ con
  $\E Z^2 < \infty$ y $0 < \theta <
        1$,
        \[
        \P\bigl(Z > \theta\,\E Z\bigr) \;\geq\; (1 -
        \theta)^2\,\frac{(\E Z)^2}{\E Z^2}
        \]
        \emph{(pártase $\E Z$ en el nivel $\theta\E Z$ y aplíquese
        Cauchy--Schwarz al trozo superior)}. (b) Demuéstrese
        $\E T_n^4 \leq 3s_n^4$. (c) Dedúzcase
        $\P\bigl(\abs{T_n} > \frac{s_n}2\bigr)
        \geq \frac3{16}$ y conclúyase que $\sum_nx_n\varepsilon_n$
        diverge c.s.; de ahí la dicotomía
        \[
        \sum_nx_n\varepsilon_n\ \text{converge c.s.}
        \iff \sum_nx_n^2 < \infty .
        \]
  \item (Serie armónica aleatoria) Conclúyase que
  $\sum_n\frac{\varepsilon_n}{n^s}$ converge c.s.\ si y solo si
  $s > \frac12$. Para $\frac12 < s \leq
        1$ la serie converge c.s.\ aunque $\sum_nn^{-s} = \infty$: los
  signos aleatorios producen cancelación de intensidad raíz cuadrada ---
  compárese con la serie alternada $\sum_n\frac{(-1)^n}{n^s}$, que
  converge para \emph{todo} $s > 0$.
\end{enumerate}

\medskip
\noindent\textbf{Parte VII --- Concentración: la desigualdad de
Hoeffding.} La \omterm{def:b3:probability:space}{ley} fuerte dice
$\frac{S_n}n \to m$; las desigualdades de concentración dicen cuán
improbable es una desviación \emph{para cada $n$ fijo}.
\begin{enumerate}[resume]
  \item (Lema de Hoeffding) (a) Demuéstrese
  $\cosh\lambda \leq \eu^{\lambda^2/2}$ para todo $\lambda \in \R$,
  comparando las dos series término a término. (b) Sea $Z$ centrada con
  $a \leq Z \leq b$, $a < b$. Demuéstrese
        \[
        \E\,\eu^{\lambda Z} \leq
        \exp\Bigl(\frac{\lambda^2(b - a)^2}8\Bigr)
        \]
        \emph{(acótese $\eu^{\lambda z}$ en $\intcc ab$ por su cuerda,
        tómense \omterm{def:b3:probability:space}{esperanzas} y
        estúdiese $\varphi(t)
        = -pt + \log(1 - p + p\eu^t)$ con $p =
        \frac{-a}{b-a}$ y $t = \lambda(b - a)$: demuéstrense
        $\varphi(0) = \varphi'(0) = 0$ y $\varphi''
        \leq \frac14$)}. \item (Desigualdad de Hoeffding) Sean
        $X_1, \dots, X_n$
        \omterm{def:b3:probability:independence}{independientes} con
        $a_i \leq X_i \leq b_i$ y $S_n =
        X_1 + \dots + X_n$. Demuéstrese, para $t > 0$,
        \[
        \P\bigl(S_n - \E S_n \geq t\bigr) \leq
        \exp\Bigl(\frac{-2t^2}{\sum_{i=1}^n(b_i -
        a_i)^2}\Bigr),
        \]
        y la misma cota para la cola inferior \emph{(Chebyshev
        exponencial: acótese $\E\,\eu^{\lambda(S_n - \E S_n)}$ usando la
        \omterm{def:b3:probability:independence}{independencia} y la
        pregunta 17, y optimícese después en
        $\lambda > 0$)}.
  \item (La \omterm{def:b3:probability:space}{ley} fuerte, caso acotado,
  con velocidad) Sean las $X_i$ i.i.d.\ con valores en $\intcc ab$ y $m
        = \E X_1$. Demuéstrese
        \[
        \P\Bigl(\Bigl|\frac{S_n}n - m\Bigr| \geq
        \varepsilon\Bigr) \leq
        2\exp\Bigl(\frac{-2n\varepsilon^2}{(b - a)^2}\Bigr)
        \]
        y recupérese $\frac{S_n}n \to m$ c.s.\ por Borel--Cantelli: una
        segunda demostración de la
        \omterm{def:b3:probability:space}{ley} fuerte para variables
        acotadas --- sin truncamiento, con una velocidad exponencial en
        cada $n$ finito, pero con sumandos acotados y
        \omterm{def:b3:probability:independence}{independencia}
        \omterm{def:b3:complete:complete}{completa}. Compárense las hipótesis con las de Etemadi. \item
        (\omterm{ex:b3:probability:sllnapps}{Monte Carlo}, garantizado a
        $n$ fijo) Sean $g
        \colon \intcc01^d \to \intcc01$
        \omterm{def:b3:lebesgue:measurable}{medible} y $(U_k)$ la
        muestra uniforme i.i.d.\ de la pregunta 11. Dado
        $\varepsilon, \delta > 0$, demuéstrese
        \[
        n \geq \frac{\log(2/\delta)}{2\varepsilon^2}
        \implies
        \P\Bigl(\Bigl|\frac1n\sum_{k=1}^ng(U_k) -
        \int g\,\dd\lambda_d\Bigr| \geq \varepsilon\Bigr)
        \leq \delta,
        \]
        y evalúese el umbral para $\varepsilon =
        \delta = 10^{-2}$. La cota no involucra $d$: compárese con la
        pregunta 11 y con las mallas deterministas.
\end{enumerate}

\medskip
\noindent\textbf{Parte VIII --- ¿Cuán grande es un paseo aleatorio?
Hacia el logaritmo iterado.} Sea $S_n = \varepsilon_1 +
\dots + \varepsilon_n$ el paseo aleatorio simple construido con signos
equilibrados i.i.d.
\begin{enumerate}[resume]
  \item (Colas subgaussianas) Demuéstrese $\E\,\eu^{\lambda S_n} =
        (\cosh\lambda)^n \leq \eu^{n\lambda^2/2}$ y dedúzcase, para
  $x > 0$,
        \[
        \P(S_n \geq x) \leq \eu^{-x^2/(2n)},
        \qquad
        \P(\abs{S_n} \geq x) \leq 2\,\eu^{-x^2/(2n)} .
        \]
  \item Dedúzcase, vía Borel--Cantelli,
        \[
        \limsup_{n\to\infty}\frac{\abs{S_n}}
        {\sqrt{2n\log n}} \leq 1 \quad\text{c.s.}
        \]
        \emph{(para $\eta > 0$, súmense las cotas de cola en $x =
        (1 + \eta)\sqrt{2n\log n}$ e interséquese después en
        $\eta = \frac1p$)}. En particular, el paseo vive a la escala del
        TCL $\sqrt n$ salvo un factor logarítmico --- muy por debajo de
        la cota burda $\abs{S_n} \leq n$. \item A lo largo de la
        subsucesión de duplicación $n_j = 2^j$, demuéstrese
        \[
        \limsup_{j\to\infty}\frac{S_{n_j}}
        {\sqrt{2n_j\log\log n_j}} \leq 1 \quad\text{c.s.},
        \]
        y reflexiónese: la \emph{\omterm{def:b3:probability:space}{ley}
        del logaritmo iterado} (Khinchin; Hartman--Wintner para sumandos
        $L^2$ centrados generales) afirma que
        \[
        \limsup_{n\to\infty}\frac{S_n}
        {\sqrt{2n\log\log n}} = 1 \quad\text{c.s.}
        \]
        Explíquese con precisión qué separa la estimación por
        subsucesiones recién demostrada de la mitad superior de este
        enunciado (hay que controlar $\max_{n_j \leq n \leq
        n_{j+1}}S_n$ dentro de cada bloque, lo que exige una desigualdad
        maximal a escala \emph{exponencial}) y compruébese
        cuantitativamente que la desigualdad de la pregunta 12 es
        demasiado débil para ese fin. La mitad inferior descansa en el
        segundo lema de Borel--Cantelli aplicado a bloques
        \omterm{def:b3:probability:independence}{independientes}; ambas
        mitades son material honesto de tercer año para un curso
        dedicado a la probabilidad.

  \item (Desviación uniforme sobre una clase finita) Sean $A_1,
        \dots, A_N$ sucesos de un experimento repetible, y estímese cada
  probabilidad por su frecuencia empírica $\hat p_i$ sobre $n$
  repeticiones i.i.d. Combinando la desigualdad de Hoeffding con una
  cota de la unión, demuéstrese
        \[
        \P\Bigl(\max_{i\leq N}\,\abs{\hat p_i - \P(A_i)} >
        \varepsilon\Bigr) \;\leq\; 2N\,\eu^{-2n\varepsilon^2},
        \]
        y dedúzcase la regla del tamaño muestral: $n \geq
        \frac{\ln(2N/\delta)}{2\varepsilon^2}$ garantiza que las $N$
        estimaciones sean simultáneamente $\varepsilon$-precisas con
        probabilidad $\geq 1 -
        \delta$. Calcúlese $n$ para $N = 10^6$, $\varepsilon =
        0.01$, $\delta = 0.05$: el precio logarítmico de la uniformidad.
        \item (La ventana armónica aleatoria) Combinando las dos mitades
        de la teoría de series aleatorias, demuéstrese que, para signos
        i.i.d.\ $(\varepsilon_n)$, la serie
        $\sum_n\frac{\varepsilon_n}{n^\alpha}$ converge c.s.\ si
        $\alpha > \frac12$ y diverge c.s.\ si $\alpha \leq \frac12$;
        contrástese con la convergencia absoluta (que exige
        $\alpha > 1$): en la ventana $\alpha \in \intoc{\frac12}1$, la
        convergencia es un fenómeno genuinamente probabilístico ---
        cancelación, no tamaño.
\end{enumerate}
\end{problem}
